% LaTeX companion of 11-eigenvalues-and-eigenspaces.typ. Both formats are editable.
\chapter{Eigenvalues and eigenspaces}

\section{WS 24, p. 1: characteristic polynomial and multiplicities}

\begin{theorem}{Eigenvalues are roots}

The eigenvalues of \(T\) are precisely the roots of its characteristic polynomial.

\end{theorem}

简明 proof: \(\lambda\) 是 \(T\) 的 eigenvalue iff \(E_{\lambda} \neq \left\{ 0 \right\}\) iff \(\text{ker}\left( {T - \lambda I} \right) \neq \left\{ 0 \right\}\) iff nullity\(\left( {T - \lambda I} \right) > 0\) iff \(\text{rank}\left( {T - \lambda I} \right) < n\) iff \(T - \lambda I\) 不可逆 iff \(\text{det}\left( {T - \lambda I} \right) = 0\). Thus \(\lambda\) is a root of \(\text{det}\left( {T - \lambda I} \right) = 0\).

\begin{theorem}{Distinct-eigenvalue eigenspaces}

\(\forall x \in E_{\lambda_{1}}\) 且 \(x \neq 0\), and \(y \in E_{\lambda_{2}}\) 且 \(y \neq 0\), if \(\lambda_{1} \neq \lambda_{2}\), then \(x,y\) are linearly independent.

\end{theorem}

因而 eigenbases of distinct eigenspaces have linearly independent union, and if \(\sum_{i}\text{dim}\left( E_{\lambda_{i}} \right) = \text{dim}V\), that union is an eigenbasis of \(V\) for \(T\).

\begin{theorem}{Geometric and algebraic multiplicity}

For eigenvalue \(\lambda\), \(\text{gem}(\lambda) \leq \text{alm}(\lambda)\).

\end{theorem}

Let \(A \in \mathbb{R}^{n \times n}\), let \(\lambda\) be an eigenvalue and \(\text{gem}(\lambda) = m\). Let \(\left( {v_{1},\ldots,v_{m}} \right)\) be a basis of \(E_{\lambda}\), put \(S = \begin{pmatrix}
v_{1} & \ldots & v_{m}
\end{pmatrix}\), and define \(B = S^{- 1}AS\). Then for \(1 \leq i \leq m\),

\[Be_{i} = S^{- 1}ASe_{i} = S^{- 1}Av_{i} = S^{- 1}\lambda v_{i} = \lambda e_{i}.\]

Therefore

\[B = \begin{pmatrix}
{\lambda I_{m}} & P \\
0 & Q
\end{pmatrix},\]

so

\[f_{A{(x)}} = f_{B{(x)}} = \text{det}\left( {B - xI_{n}} \right) = \left( {x - \lambda} \right)^{m}f_{Q{(x)}},\]

and \(\text{alm}(\lambda) \geq m\).

\begin{theorem}{Degree of the characteristic polynomial}

For \(T:V\rightarrow V\), \(\chi_{T{(x)}}\) has degree \(\text{dim}V\).

\end{theorem}

This is the corollary of the preceding two theorems: if \(T\) has \(n\) different eigenvalues, then \(\sum_{i}\text{alm}\left( \lambda_{i} \right) \leq \text{dim}V\). Proof: 见 P9.
